\chapter{Discovery}
\label{ch:discovery}

\section*{What this chapter is for}

Every build round starts with a brief that is too vague to build from, and that
vagueness is deliberate. Turning it into something specified is the first thing
scored and the thing most candidates rush past on their way to the part they
are comfortable with.

\section{Why the brief is vague on purpose}

A brief that could be implemented directly would test implementation. The gap
between what you are told and what you would need to know is the exercise, and
it maps directly onto the level: Chapter~\ref{ch:staff} defined Staff by
unscoped problems, and this is an unscoped problem with a clock on it.

\judgement{An engineer who starts building in the first five minutes has
answered the question ``can you implement?'' when the question was ``can you
decide what to implement?''. The work may be excellent and it is evidence for
the wrong thing.}

\section{The trap: do not arrive with a data model}

The strongest instinct, once you know roughly what domain the exercise is in,
is to prepare the schema. It is a few hours' work, it removes the riskiest part
of the session, and it is the single most expensive preparation mistake
available.

\begin{kittrap}
A prepared data model destroys the first scored criterion. Requirements
discovery cannot be demonstrated by somebody who arrives with the requirements
already discovered, and producing a complete schema in minute three does not
read as preparedness \dash{} it reads as a candidate who has stopped listening.

Worse, it makes you brittle. A prepared model is a model of the problem you
guessed at, and the brief will differ. Then you are either defending a design
against a requirement it does not fit, or visibly abandoning it.
\end{kittrap}

\begin{kitmethod}
\textbf{What you may prepare:} the method \dash{} what to ask and in what
order. The vocabulary of the domain. Your tooling, project skeleton, and any
standards you work to. A question list. A run sheet. Every one of these is
about \emph{how} you work and exists independently of this exercise.

\textbf{What you may not prepare:} the entities, the relationships, the
endpoints, the acceptance criteria. Those are the answer, and they are what the
session is for.
\end{kitmethod}

\section{The questions, in order}

The order matters, because the early ones constrain the later ones and asking
them in the wrong order produces answers you have to discard.

\begin{kitmethod}
\textbf{1. Who is this for, and what do they do today?} Everything else depends
on it, and the answer is frequently surprising.

\textbf{2. What is the one thing that has to work?} Forces a priority out of
somebody who presented a list, and licenses everything you cut later.

\textbf{3. What happens when it goes wrong?} The first question that produces
domain knowledge rather than a restatement of the brief.

\textbf{4. How many, and how often?} The scale question. The answer is often
``not many'', which is licence to stop designing for a scale nobody has.

\textbf{5. Who else touches this?} Finds the second actor, who is almost always
in the brief implicitly and never named.

\textbf{6. What is deliberately not in scope?} Get it said out loud. It is the
thing you point back at when you cut something later.
\end{kitmethod}

Two more that are worth asking once each, in the domain this book is about:
what happens when the model gets it wrong, and how anybody would know. If the
brief involves a model, that pair separates a candidate who has operated one
from a candidate who has built one.

\section{What discovery has to produce}

Not notes. A short written artefact, on the shared screen, that you and they
can both point at: what is being built, for whom, the three or four acceptance
criteria, and what is explicitly out of scope.

\judgement{The out-of-scope list is the most valuable line in it, and it is the
one people leave out. Written down at minute twenty, it is a decision you made
together. Invoked at minute eighty without having been written down, it is an
excuse.}

Ten minutes on this, in a session of two hours, is proportionate. It also gives
the review at the end something to be graded against that is not just the
running code.

\section{Closing discovery}

Discovery ends. It does not taper.

Say so explicitly \dash{} ``I have enough to start; I am going to build X
first, and I will come back to Y'' \dash{} and then start. Two failure modes
sit either side of that sentence: stopping too early, and building the wrong
thing confidently; and never stopping, which produces a beautifully understood
problem and nothing running.

\judgement{Of the two, the second is more common in senior candidates and more
damaging, because the round has an artefact in its criteria and interest in the
problem is not a substitute for one.}

\begin{drill}
Take any one-sentence product brief \dash{} a friend's business, something you
use \dash{} and run the six questions against somebody playing the owner, in
ten minutes, ending with the written artefact. Then check: did you get a second
actor out of them? Did you get an explicit out-of-scope line? Those two are the
ones that go missing under time pressure.
\end{drill}

\section*{What this chapter does not claim}

That requirements discovery is scored is reported for the loops that publish
their criteria; the six questions and their order are the author's method, not
a measured optimum.

The instruction not to prepare a data model is a judgement, and a strong one.
It follows from what the round says it is measuring rather than from any
observation of candidates who did it.
